\subsection{Event Definition}\label{sec:EVENT_DEFINITION}
We define a low-energy event as the combination of a \texttt{primary cluster} (a 3D-cluster, after the plane matching of~\cref{sec:CLUSTERING_TPC}) together with its \texttt{adjacent clusters} and a list of clustered \gls{pds} hits (\texttt{flashes}) that are matching candidates for the drift coordinate reconstruction (see \cref{sec:FLASH_MATCHING}); from this list a single flash is selected as the best match. In any neutrino analysis it will become important to try to distinguish (at least on average) the \texttt{adjacent cluster} distributions for any low-energy event, not only to distinguish background activity (e.g. neutron captures) from signal events, but also to identify the different interaction channels of the neutrino interaction (\gls{cc} events produce de-excitation gammas that are not present in \gls{es} interactions). Consequently, the higher the projected background activity in any of the isolated drift volumes of \gls{dune}'s \gls{fd}s, the harder particle identification and channel separation will become. Notice that there is nothing preventing a single neutrino interaction from producing more than one low-energy event. Rarely, secondary gammas scatter further than \SI{1}{\meter} in \gls{lar} or a background-induced mismatch in the \gls{tpc} reconstruction leads to hits appearing in uncorrelated regions of the detector. Independently of this possibility, a low-energy event reconstruction efficiency can be defined as the possibility of a single particle interaction to produce at least one low-energy event that can be fully reconstructed. Here, the preselection efficiency is defined as the probability of generating at least one reconstructed cluster (after the plane matching of~\cref{sec:CLUSTERING_TPC}) with a multiplicity larger than a specific threshold. This efficiency is shown in \cref{fig:DETECTION_EFFICIENCY} for both \gls{hd} and \gls{vd} geometries as a function of the true electron energy, for a scan of the cluster's \texttt{nhit} threshold value. The chosen signal clustering parameters are 3 adjacent channels and \SI{12.5}{\micro\second} time window. The results show a clear and expected decrease in efficiency with increasing \texttt{nhit} threshold, as low-energy events are more likely to be lost when applying stricter cuts on the minimum number of hits per cluster. This shows the importance of a simple clustering approach to low-energy neutrino reconstruction, as more sophisticated algorithms might lead to a considerable decrease in efficiency. Comparing between detector geometries, one notices a slight advantage of the \gls{vd} module, caused by a lower threshold of the readout and a simpler induction-collection plane matching.

\begin{figure}[ht!]
    \centering
    \includegraphics[width=0.49\linewidth]{Images/LowE/preselection_efficiency_hd_1x2x6_marley_official_values_efficiency_electronk_nhits_tree_reco_scan.png}
    \includegraphics[width=0.49\linewidth]{Images/LowE/preselection_efficiency_vd_1x8x14_3view_30deg_marley_official_values_efficiency_electronk_nhits_tree_reco_scan.png}
    \caption{\gls{hd} (left) and \gls{vd} (right) detection efficiency as a function of the true electron energy, for signal clustering parameters of 3 adjacent channels and \SI{12.5}{\micro\second}. One curve is shown per minimum-hit threshold.\label{fig:DETECTION_EFFICIENCY}}
\end{figure}

The first figure of merit to be extracted from the simulation is the detection efficiency (see~\cref{fig:DETECTION_EFFICIENCY}). We observe a strong dependence on the number of hits. With a minimum hit condition of 1 hit, the efficiency exceeds \SI{90}{\percent} for true electron energies above \SI{3}{\mega\electronvolt} in \gls{hd} and above \SI{1}{\mega\electronvolt} in \gls{vd}; with 3 hits it reaches \SI{90}{\percent} only above about \SI{12}{\mega\electronvolt}. With the 1-hit condition, the efficiency rises from zero to its plateau over the first few \si{\mega\electronvolt} of electron energy: the effective detection threshold, the electron energy at which the efficiency reaches about \SI{50}{\percent}, is about \SI{1}{\mega\electronvolt} in \gls{hd} and about \SI{0.5}{\mega\electronvolt} in \gls{vd}. Both lie well below the \SI{4}{\mega\electronvolt} \gls{dune} threshold, which is a neutrino-energy threshold set by the radiological background rather than by the detector response, so that the reconstruction is fully efficient over the neutrino energies targeted by the analyses.

\begin{figure}[ht!]
    \centering
    \includegraphics[width=.49\linewidth]{Images/Reco/clustering_efficiency_energy_hd_1x2x6_centralapa_hd_1x2x6_lateralapa_vd_1x8x14_3view_30deg_nominal_marley_final_background_completeness_scan.png}
    \includegraphics[width=.49\linewidth]{Images/Reco/clustering_efficiency_energy_hd_1x2x6_centralapa_hd_1x2x6_lateralapa_vd_1x8x14_3view_30deg_nominal_marley_final_purity_background_purity_scan.png}
    \caption{Clustering completeness (left) and purity (right) of the primary cluster for the signal plus background samples, as a function of the true neutrino energy, for the \gls{hd} Central, \gls{hd} Lateral and \gls{vd} Nominal configurations.\label{fig:EVENT_COMPLETENESS}\label{fig:EVENT_PURITY}}
\end{figure}

The next parameters to evaluate are the completeness and purity of the main clusters. The completeness is a measure of the electron charge--confinement given a certain clustering parameter set: the fraction of the electron charge collected in the primary cluster. The purity is the fraction of the cluster charge that originates from the electron, a measure of the background contamination of the cluster. For the purpose of the solar neutrino analysis of \gls{dune}, clustering parameters of 3 adjacent channels and 25 ticks (\SI{12.5}{\micro\second}) have been used to maximise completeness while ensuring cluster purity in the context of background. The most up-to-date values for this clustering should be found in the producer's \href{https://github.com/DUNE/dunereco/blob/develop/dunereco/LowEUtils/fcl/SolarCluster.fcl}{SolarCluster.fcl}. With the full background model (\cref{fig:EVENT_COMPLETENESS}), the completeness is \SIrange{94}{96}{\percent} in \gls{hd} and \SIrange{75}{89}{\percent} in \gls{vd}, rising with energy, and the charge purity stays above \SI{99.8}{\percent} in \gls{hd} and \SI{99.5}{\percent} in \gls{vd} at all neutrino energies and drift distances. Despite the good overall completeness of the signal, we observe a trend with energy (and hence with the number of hits per cluster), which can be accounted for in the later described charge correction factor to improve reconstruction. The trend is clear and understandable: for depositions which result in a lower number of hits (this can happen because of low energy deposition or a perpendicular trajectory of the particle with respect to the \gls{tpc} plane) it is more likely that clusters of similar sizes are formed (e.g. 2 clusters of 1 hit) resulting in a lower charge in the main cluster. Both the completeness and the purity stay nearly constant along the drift, so the pile-up with background activity accumulated over the drift does not degrade the clustering at the simulated rates.

\gls{dune}'s solar analysis workflow starts with calibrating the clustered hits from the \gls{cc} interaction of simulated neutrino events with the \gls{marley} generator. The \gls{cc} channel is chosen to be the default based on its dominant cross-section in the low-energy range, as well as its richer topology. All energy and spacial resolution results presented here should be considered an underestimation (or at least compatible) to the ones results for the \gls{es} channel. Each cluster's information is reduced to a point-like object characterized by its number of hits (\texttt{nhits}), \texttt{charge} (summed over the charge of the individual hits) and a 2.5D-reconstructed vertex $<t, y, z>$ computed from the charge-weighted average of the individual hits.

For the purpose of this work, secondary low-energy events coming from a single neutrino interaction are identified using truth-level information and ignored in favour of the one with the highest charge in the primary cluster. This is to ensure that the reconstruction performance is evaluated for the most complete reconstruction object in any interaction but in any physics analysis these should be kept and treated as part of the signal or even background samples.

\begin{table}[ht!]
    \centering
    Adjacent Cluster Distribution - Vertex Distance \SI{20}{\centi\meter}\\
    \vspace{.15cm}
    \begin{tabular}{lccc}
\toprule
Configuration & \multicolumn{3}{c}{Source} \\
 & Signal & Intrinsic & External \\
\midrule
\emph{HD Central} & \num{0.5458(15)} & \num{0.0182(3)} & \num{0.0134(2)} \\
\emph{HD Lateral} & \num{0.6083(16)} & \num{0.0166(3)} & \num{0.0229(3)} \\
\emph{VD Nominal} & \num{0.5031(15)} & \num{0.1473(8)} & \num{0.2238(10)} \\
\emph{VD Shielded} & \num{0.533(2)} & \num{0.1446(11)} & \num{0.0526(7)} \\
\bottomrule
\end{tabular}

    \caption{Average number of adjacent clusters per primary signal cluster for \SI{20}{\centi\meter} distance, split by origin (signal de-excitation products, intrinsic and external background).\label{tab:ADJCLUSTERNUM_RADIAL_20}} 
\end{table}

For the specific purpose of the solar neutrino studies an analyzer module (\href{https://github.com/DUNE/duneana/blob/develop/duneana/SolarNuAna/SolarNuAna_module.cc}{SolarNuAna\_module.cc}) has been written to combine these reconstruction objects with the truth information of the event. This module essentially combines all the information needed to perform the calibration and reconstruction of the signal candidates. The biggest caveat from the analysis point of view is that all clusters and flash information is reduced to a point-like data structure. This is, of course, a choice of design in the analyzer and allows for a simpler handling of the data because all event candidates have the same dimension. In future works, this analyzer should be substituted by the more general \href{https://github.com/DUNE/duneana/blob/develop/duneana/LowEAna/LowEAna_module.cc}{LowEAna\_module.cc} which aims to provide the same event reconstruction while keeping the internal structure of the candidate clusters to allow for better calibration and discrimination strategies.

\begin{figure}[ht!]
    \centering
    \includegraphics[width=0.9\linewidth]{Images/Reco/hd_1x2x6_centralapa_marley_edep_edep_tpc_event_2_display_hd_1x2x6_centralapa_marley_edep_z_y_x_pdg_event_display.png}
    \caption{Example of a simulated neutrino interaction in the \gls{hd} geometry, in two \gls{tpc} projections, showing the primary electron (orange) and the associated de-excitation gammas (blue). True energy deposits are marked with filled circles and reconstructed clusters with empty circles.\label{fig:EXAMPLE_EVENT}}
\end{figure}

An example of a simulated neutrino interaction in the \gls{hd} geometry is shown in~\cref{fig:EXAMPLE_EVENT}, now showing the reconstructed clusters. Compared to~\cref{fig:EVENT_EDEP}, the clusters are reduced to a point-like data structure, which results from the clustering algorithm and aligns with the high density regions of the simulated deposits.
